% Complete proofs for the separately attributed modern finite-algebra companion.
\section{Finite-Algebra and Context-Congruence Proofs}
\label{sec:p09-tree-algebra-annex}

\begin{annexitem}{annex:p09.tree-algebra.evaluation}{thm:p09-tree-algebra-evaluation}{Ground-term evaluation and generated carrier}
\begin{proof}
Define the value of a nullary term \(f\) to be \(f^A\).
Having defined the values of all children, define the value of
\(f(t_1,\ldots,t_n)\) by the displayed operation. Ground terms are
finite well-founded syntax trees, so recursion on their height defines
a value for every term. Each operation is total on the given carrier,
so every recursive application is well defined. These definitions
preserve the term-forming operations, hence define a homomorphism.

For uniqueness let \(h:T(F)\to A\) be any homomorphism.
It must send a constant \(f\) to \(f^A\). If its values on the
children equal the recursively defined ones, operation preservation
forces the same value at their parent. Induction on height therefore
gives \(h=\varphi_A\).

The image contains every constant value. If \(r_i=\varphi_A(t_i)\)
are image elements, then \(f^A(r_1,\ldots,r_n)\) is the image of
\(f(t_1,\ldots,t_n)\), so the image is closed under every operation.
Conversely any subalgebra containing the constant values contains
\(\varphi_A(t)\) for every term, by the same induction.
It is therefore the generated subalgebra. Since \(F_0\ne\varnothing\),
at least one constant witnesses its nonemptiness.

For the finite computation begin with all constant values and repeatedly
adjoin operation values on already obtained tuples. At most
\(\operatorname{card}(|A|)\) distinct values can be added.
An initial value has its constant witness; an added value has the term
formed from the previously recorded witnesses. Conversely induction
on each term shows that its value eventually appears. When the closure
stabilizes it is exactly the image, and all its elements have the
specified finite witnesses.
\end{proof}
\end{annexitem}

\begin{annexitem}{annex:p09.tree-algebra.contexts}{thm:p09-tree-congruence-contexts}{Context stability and the coarsest saturating congruence}
\begin{proof}
First suppose \(\sim\) is an operator congruence and \(u\sim v\).
Induct on the context. The trivial context preserves the relation.
A nontrivial context has the form
\(f(t_1,\ldots,t_{i-1},D,t_{i+1},\ldots,t_n)\), where
\(D\) has one hole and the other children are ground terms.
By induction \(D[u]\sim D[v]\); reflexivity relates the equal
other children. Operator compatibility gives \(C[u]\sim C[v]\).
This proves stability under every one-hole context.

Conversely suppose an equivalence relation is stable under all such
contexts. Given \(u_i\sim v_i\), replace the children one at a time.
At position \(i\) use the context
\(f(v_1,\ldots,v_{i-1},x,u_{i+1},\ldots,u_n)\).
Context stability relates the terms before and after this replacement.
Transitivity through all \(n\) steps relates the two full terms.
For a nullary symbol the required assertion is reflexivity.
Thus the relation is an operator congruence.

For \(\equiv_L\), reflexivity, symmetry and transitivity follow from
the corresponding properties of equality of truth values
\(C[t]\in L\), for each fixed \(C\). If \(u\equiv_L v\), then,
for any \(D,C\), the composed one-hole context \(C[D]\) gives
\[
 C[D[u]]\in L\quad\Longleftrightarrow\quad C[D[v]]\in L.
\]
Thus \(D[u]\equiv_L D[v]\). The first part makes it a congruence.
The trivial context gives \(u\in L\) iff \(v\in L\), so it
saturates \(L\).

Finally, if a congruence \(\sim\) saturates \(L\) and \(u\sim v\),
context stability gives \(C[u]\sim C[v]\) for every \(C\).
Saturation implies agreement of their membership in \(L\).
Hence \(u\equiv_L v\), proving \(\sim\subseteq\equiv_L\).
\end{proof}
\end{annexitem}

\begin{annexitem}{annex:p09.tree-algebra.recognition}{thm:p09-tree-algebra-recognition}{All four recognition conditions}
\begin{proof}
From a finite bottom-up recognizer use the complete subset construction
in \cref{sec:p09-finite-tree-determinization}. Its carrier is the
nonempty finite set \(\mathcal P(Q)\), including \(\varnothing\);
its total \(f\)-operation takes child reachable-state subsets to
the subset of all permitted target states. Its designated subset
consists of the subsets meeting the old final set. The proved subset
invariant identifies its evaluation with the possible old root states,
so it recognizes exactly the old language. This proves (1) implies (2)
without assuming the original machine was deterministic or complete.

Given (2), define \(u\sim v\) iff \(\varphi(u)=\varphi(v)\).
It is an equivalence relation. If corresponding child values agree,
applying the same algebra operation gives equal parent values.
Thus it is a congruence. Its classes correspond to elements of the
finite image, so its index is finite. The inverse-image condition
makes \(L\) a union of its classes. This proves (2) implies (3).

For (3) let \(\sim\) be the saturating finite-index congruence.
By the preceding lemma it refines \(\equiv_L\). The map
\([t]_\sim\mapsto[t]_{\equiv_L}\) is well defined and surjective.
The latter set of classes is therefore finite, proving (4).

For (4), the preceding lemma proves that \(\equiv_L\) is a
congruence and saturates \(L\). Its quotient has a nonempty finite
carrier, since \(T(F)\) contains a constant. Given representatives
\(t_i,u_i\) of the same classes, compatibility gives
\([f(t_1,\ldots,t_n)]=[f(u_1,\ldots,u_n)]\); hence its displayed
operations are well defined and total. Saturation makes its
designated subset well defined. The map \(t\mapsto[t]\) preserves
every operation, and membership in its designated subset is exactly
\(t\in L\). This proves (4) implies (2) and verifies the specified
canonical quotient.

Finally, a finite recognition algebra supplies a deterministic complete
machine with its carrier as states, operation tables as transitions,
and \(A'\) final. At a leaf its unique state is its constant value;
induction on the children gives a unique state at every parent,
namely its evaluation value. Its root is final iff its evaluated
element is in \(A'\). This proves (2) implies (1) and closes all
directions.
\end{proof}
\end{annexitem}

\begin{annexitem}{annex:p09.tree-algebra.minimal}{thm:p09-tree-algebra-minimal}{Canonical surjection and uniqueness}
\begin{proof}
Let \(R=\varphi(T(F))\). The evaluation lemma makes it a nonempty
subalgebra. The kernel of \(\varphi\) is a congruence saturating \(L\),
as just proved, so it refines \(\equiv_L\). Consequently
\(\varphi(t)=\varphi(u)\) implies \([t]=[u]\), making the
formula \(\pi(\varphi(t))=[t]\) independent of the chosen witness.
It is surjective because every quotient class has a representative.
It is unique because every element of \(R\) is a ground-term value.

For \(r_i=\varphi(t_i)\), operation preservation and the quotient
definition give
\[
 \pi(f^A(r_1,\ldots,r_n))
   =\pi(\varphi(f(t_1,\ldots,t_n)))
   =[f(t_1,\ldots,t_n)]
   =f^{A_L}(\pi(r_1),\ldots,\pi(r_n)).
\]
The nullary case gives the same assertion for constants.
Thus \(\pi\) is a homomorphism. For \(r=\varphi(t)\),
\[
 r\in A'\cap R
 \quad\Longleftrightarrow\quad t\in L
 \quad\Longleftrightarrow\quad \pi(r)\in Q_{\min,f}.
\]
This proves exact preservation of recognition.

A surjection between finite sets gives
\(\operatorname{card}(R)\geq\operatorname{card}(Q_{\min})\).
If equality holds it is a bijection. Its inverse preserves operations:
apply \(\pi\) to an operation on inverse images, use the displayed
identity, then use injectivity. The membership identity likewise
preserves the designated set in both directions. Therefore it is
an isomorphism of recognition algebras.

Unreachable states of any complete deterministic recognizer can be
removed: the ground-term image is closed under every operation and
contains every constant, so restricting tables to it remains total
and preserves all ground-term evaluations. A minimum complete recognizer
therefore has no unreachable state. The quotient recognizer is
reachable, since every \([t]\) is its evaluation of \(t\).
It attains the lower bound, and every recognizer attaining that bound
is isomorphic to it. This proves minimality and uniqueness up to
state renaming in precisely the complete deterministic class.
\end{proof}
\end{annexitem}

\begin{annexitem}{annex:p09.tree-algebra.refinement}{thm:p09-tree-algebra-refinement}{Full refinement correctness}
\begin{proof}
The input is complete, finite and reduced: every state is the evaluation
of some ground term. Write \(\delta(t)\) for the unique ground-term
evaluation in this state algebra. A ground context \(C\) induces a total function
\(C^\delta:Q\to Q\) by inserting a state at its hole and evaluating
the remaining operations. Recursion on \(C\) gives
\(C^\delta(\delta(t))=\delta(C[t])\).

The base \(P_0\) relates states with the same final/nonfinal status.
It is an equivalence relation, and the trivial context is exactly
the hole-depth-zero test, so the asserted invariant holds for \(j=0\).

Assume the invariant for \(P_j\). An elementary one-hole operation
is the function
\[
 \tau(q)=\delta_f(q_1,\ldots,q_{i-1},q,q_{i+1},\ldots,q_n).
\]
Its other states have ground-term witnesses by reducedness. It is
therefore induced by a ground context whose hole depth is one.
The recurrence for \(P_{j+1}\) is the intersection of \(P_j\)
and all inverse-image relations
\(\{(q,r):\tau(q)P_j\tau(r)\}\).
Each is an equivalence relation, so their intersection is one.
Only finitely many such functions need testing because the signature
and state set are finite.

If \(qP_{j+1}r\), the old relation handles contexts of depth at
most \(j\). Any nontrivial context of depth at most \(j+1\) decomposes
as \(D[E]\), where \(E\) is the operation at the parent of the
hole with its other ground children fixed, and \(D\) has hole depth
at most \(j\). The recurrence gives
\(E^\delta(q)P_jE^\delta(r)\), and the induction hypothesis for
\(D\) yields equal final status of the complete evaluations.
Thus every required context agrees.

Conversely suppose all contexts of depth at most \(j+1\) agree.
They include every context of depth at most \(j\), so \(qP_jr\).
For each tested \(\tau\) and every ground context \(D\) of depth
at most \(j\), the composition of \(D\) with the witnessing
one-hole context for \(\tau\) has depth at most \(j+1\).
Agreement on these composites and the induction hypothesis give
\(\tau(q)P_j\tau(r)\).
All recurrence tests therefore hold. This proves the invariant at
\(j+1\) in both directions.

Every new partition refines the previous one. A strict splitting round
increases the number of nonempty blocks by at least one; it can never
exceed \(\operatorname{card}(Q)\). Starting with
\(\operatorname{card}(Q/P_0)\) blocks gives the asserted bound.
A round that makes no split is stable under the same finite test
operator, so repeating it makes no further change.

At stability \(P\), every elementary operation preserves \(P\).
Replace the inputs of an \(n\)-ary operation one coordinate at a
time; transitivity shows full operator compatibility. The nullary
case requires only reflexivity. Therefore \(P\) is a congruence
on the state algebra, quotient operations are representative independent,
and every block is final or nonfinal by refinement of \(P_0\).
Induction on terms shows that quotient evaluation is the class of the
old evaluation; recognition is consequently preserved.

Stability also says states related by \(P\) remain related in every
finite context, by induction on that context. They have the same
final status. Conversely, if states are indistinguishable by every
ground context, the proved bounded-depth invariant puts them in every
\(P_j\), hence in the stable \(P\).
For ground witnesses \(u,v\), this is exactly \(u\equiv_L v\).
Reducedness supplies witnesses for every state, so stable classes
correspond bijectively to \(Q_{\min}\), and their operations and
final flags correspond. The output is therefore the canonical minimum,
not merely an arbitrary smaller equivalent machine.
\end{proof}
\end{annexitem}

\begin{annexitem}{annex:p09.tree-algebra.cases}{thm:p09-tree-algebra-cases}{Branching calculations and the completeness countercase}
\begin{proof}
In the parity algebra, the value of leaf \(a\) is zero and that of
leaf \(b\) is one, exactly their \(b\)-leaf counts modulo two.
The count at a binary root is the sum of its child counts.
Induction on the tree therefore proves the parity evaluation.
The displayed tree has child values \(1\) and \(0+1=1\);
the root is \(1+1=0\) modulo two and is designated.
The constants reach both states. The trivial context distinguishes
\(a\) from \(b\) by acceptance, so any complete deterministic
recognizer needs at least two states. The given one has exactly two.

The same count induction gives the modulo-three and modulo-four
evaluations. In the former, \(0,1,2\) are reached respectively by
\(a,b,f(b,b)\). In the initial partition \(1,2\) are both nonfinal,
but adding state \(1\) at the right child gives states \(2,0\).
These have different final status. The first refinement splits
the nonfinal block, leaving singleton classes, which are stable.

In the modulo-four example states \(0,1,2,3\) have witnesses
\(a,b,f(b,b),f(b,f(b,b))\).
The initial blocks are \(\{0,2\}\) and \(\{1,3\}\).
Replacing a state by another of the same parity preserves the parity
of every addition modulo four, so every refinement test preserves
these two blocks. They are already stable. The quotient is the
two-state parity algebra with the given designated class.

For the leftmost algebra, a constant has its own label type.
At an \(f\)-root, the first child contains the leftmost leaf, and
the operation selects exactly its type. Structural induction proves
the assertion for all terms. The two displayed roots consequently
evaluate to \(0\) and \(1\), respectively, and have different acceptance.

Finally, over the nullary alphabet \(\{a,b\}\), the only ground terms
are those two constants. The partial machine's only transition reaches
its final state on \(a\), while \(b\) has no run. Thus its language
is \(\{a\}\). Zero states could supply no accepting run on \(a\),
so its one state is minimal among partial machines.
A complete deterministic machine must assign a state to each constant.
If their states coincided, both constants would have the same final
status, contrary to the required language. It needs at least two states.
Two suffice by assigning distinct states to \(a,b\), designating only
the first. This proves the exact one-versus-two countercase.
\end{proof}
\end{annexitem}

\begin{annexitem}{annex:p09.tree-algebra.mso}{thm:p09-tree-algebra-mso}{Algebra recognition and the existing logical translations}
\begin{proof}
A finite-algebra recognition becomes a complete deterministic bottom-up
automaton by the recognition theorem. The already proved converse in
\cref{annex:p9.trees.1} describes its accepting run by existential
monadic state sets, with partition, legal transition and final-root
clauses, over exactly the stated finite-tree signature.
Hence its tree language is MSO definable.

Conversely, that annex constructs a finite bottom-up automaton for
every MSO sentence on these finite structures, using legal assignment
tracks, atomic tests, aligned Boolean operations and existential
track projection. The recognition theorem converts this automaton
to a complete finite-algebra recognizer via subset determinization.
The two directions therefore give the asserted equivalence without
a new assumption about the infinite successor tree.

For a supplied term, compute the unique algebra value by recursion
and test membership in the designated subset. For emptiness compute
the reachable subalgebra \(R\) by the evaluation lemma.
The language is nonempty iff \(R\cap A'\ne\varnothing\):
an element in this intersection has a recorded ground-term witness
accepted by the algebra, while every accepted term supplies an element
of the intersection. Both computations use finite explicit tables,
so they terminate and decide the specified problems.
All structures and witness trees here are finite. No quantification
over arbitrary infinite sets or infinite accepting runs was introduced.
\end{proof}
\end{annexitem}
